\chapter{Présentations des résultats}

\section{Essaim}

\begin{figure}[h]
    \centering
    \includegraphics[width=0.8\textwidth]{11-Résultats/pic/trajectoires_marquage.pdf}
    \caption{Essai trajectoire de l'essaim.}
    \label{fig:essaim_trajectory}
\end{figure}

\begin{itemize}
    \item afficher plot des inter distances et dire que c'est constant 
    \item montrer la trajectoire de chaque robot
    \item montrer avec une perturbation ?
\end{itemize}

\section{Essaim}

\begin{itemize}
    \item afficher plot des inter distances et dire que c'est constant 
    \item montrer la trajectoire de chaque robot
    \item montrer avec une perturbation ?
\end{itemize}

\section{Commande événementielle}

\subsection{Commande événementielle classique}

Montrer la même choque qu'avant? 
Conclusion : on diminue moins les commandes mais est ce qu'on diminue vraiment la consommation

\subsection{Commande événementielle prédictive}
Montrer la même choque qu'avant? 
Montrer la diminution des échanges

\section{Comparaisons des différentes méthodes}
